\clearpage
\chapter*{Erratum à l'exposé n\textsuperscript{o} 195}
\addcontentsline{toc}{chapter}{Erratum à l'exposé n\textsuperscript{o} 195}
\setcounter{footnote}{0}
\src{300}

\noindent\textbf{4. TDTE II : Le théorème d'existence en théorie formelle des
modules.}

\begin{center}
[Séminaire Bourbaki, t. 12, 1959/60, n\textsuperscript{o} 195, 22 p.]
\end{center}

\noindent\textbf{Page 8.} --- La formule écrite dans la proposition 5.1 n'est
correcte que lorsque $\Lambda$ est un corps ; dans le cas général, il faut
remplacer $\mathfrak m_\xi/\mathfrak m_\xi^2$ par le quotient de cet espace par
l'image de $\mathfrak n_\xi/\mathfrak n_\xi^2$, où $\mathfrak n$ est l'idéal
maximal de $\Lambda$. De plus, la définition donnée pour $O$ simple sur
$\Lambda$ n'est correcte que lorsque l'extension résiduelle $k'/k$ est
séparable. Dans le cas général, cf. SGA III, 1.1.

\noindent\textbf{Page 14, remarque.} --- Les problèmes soulevés ici sont
complètement résolus dans le cas projectif par les « schémas de Hilbert »
(cf. TDTE IV). Des exemples de Nagata et Hironaka montrent par contre que les
foncteurs envisagés ne sont pas nécessairement représentables si on ne fait
pas d'hypothèse projective, même en se bornant à la classification des
sous-variétés de dimension zéro d'une variété complète non singulière de
dimension 3 ; ceci est lié au fait que le carré symétrique d'une telle variété
peut ne pas exister.

\src{301}
\noindent\textbf{Page 15, n\textsuperscript{o} 3.} --- Pour une étude plus
complète, voir TDTE V.

\noindent\textbf{Page 16, proposition 3.1.} --- Au lieu de « si $f$ est
propre », lire « si $f$ est propre et séparable ».

\noindent\textbf{Pages 18 et 19, remarques, 2\textsuperscript{o}.} --- J'ai
montré récemment que le schéma formel des modules pour une variété abélienne
sur un corps est bien simple sur l'anneau de Witt, en d'autres termes que tout
schéma abélien sur un anneau artinien local, quotient d'un autre, provient par
réduction d'un schéma abélien sur ce dernier. La démonstration utilise
simplement les propriétés de variance de la classe d'obstruction au relèvement,
introduite dans l'exposé n\textsuperscript{o} 182, page 12. Rappelons aussi que
les schémas de modules pour les courbes de genre $g$ ou les schémas abéliens
polarisés ont été construits par Mumford (cf. SHMT).

\noindent\textbf{Page 19, paragraphe 5.} --- Il faut supposer que la section
envisagée pour définir $U$ soit non-diviseur de zéro non seulement sur
$\mathfrak X$, mais sur tout $X_n$.

\noindent\textbf{Page 20, lignes 12 et 13.} --- Au lieu de
\[
  G_x=G\bigl(\operatorname{Spec}(\widehat{\mathcal O}_{\mathfrak X,x})\bigr),
\]
lire
\[
  G_x=G\bigl(\operatorname{Spec}(\mathcal O_{\mathfrak X,x})\bigr)
  \subset
  G\bigl(\operatorname{Spec}(\widehat{\mathcal O}_{\mathfrak X,x})\bigr),
\]
et au lieu de « $\mathcal O'_{0x}$ », lire
« $\widehat{\mathcal O}'_{0x}$ ».
